\section{Introduction}
\label{sec:intro}

Vision-language-action (VLA) policies report strong success rates on standard manipulation benchmarks, yet a high aggregate success rate can hide brittle behavior under small, task-relevant changes to the scene. A policy that grasps a bowl on $74\%$ of unperturbed episodes may nonetheless fail when the same bowl is moved a few centimeters---a change a competent agent should absorb. Recent work has shown that such displacement-induced failures exist for VLA policies and framed them as a ``memory trap'' \citep{afi2025memorytrap}, showing qualitatively that the end-effector is driven toward the object's original location. We take that phenomenon as established and ask the next, \emph{causal} question---one that observing the failure cannot answer: \emph{what information actually controls the failed reach?}

A displacement failure could arise from (i) perception that mislocalizes the object, (ii) an inability to recover once the reach is underway, or (iii) a learned prior that aims at a fixed, canonical location largely independent of the current object position. These imply different remedies---better object localization, closed-loop recovery, or reshaping the action prior---yet success-rate metrics, and even the qualitative observation that the hand goes to the original location, cannot distinguish them. The distinguishing experiment is a \emph{counterfactual}: hold the physical scene fixed and change only where the object \emph{appears}.

\paragraph{Approach.} We run exactly that counterfactual. A \emph{renderer-local object-paste} intervention moves only the target object's rendered appearance to a chosen visual location while leaving its true pose, the rest of the scene, and the physics unchanged, so we can ask whether the reach follows the object's appearance or ignores it. We pair this active probe with a passive \emph{reach-field} measurement: after displacing the object by $\dv=\Tv-\Cv$ (with a fresh observation), we record the endpoint error $\ev=\Ev-\Tv$ and summarize it with $\cossim(\ev,-\dv)$ and the signed projection $\langle\ev,-\dv\rangle/\lVert\dv\rVert^2$ (near $0$ $=$ tracking the displaced object; near $1$ $=$ reaching back to the canonical location).

\paragraph{Contributions.}
\begin{itemize}
\item \textbf{Causal dissociation (main result).} Using a renderer-local object-paste counterfactual that prior work does not run, we show that a competent SmolVLA's reach endpoint is \emph{not sufficiently steered by the local rendered object position}: with the object physically canonical but shown displaced, it still grasps at the canonical location in $16/16$ trials, and the reverse paste does not rescue a displaced grasp (\cref{sec:sufficiency}). The displacement failure is located in the action prior, not in local visual object localization.
\item \textbf{Quantifying the memory trap.} We introduce a displacement-relative reach-field statistic and use it to confirm and \emph{quantify} the qualitative memory trap of \citet{afi2025memorytrap} on a competent ($74\%$) checkpoint: median $\cossim(\ev,-\dv)=0.759$ at $50$\,mm, monotonically increasing with displacement, with displaced success collapsing to $31\%$ (\cref{sec:competent}).
\item \textbf{Stability and generalization.} The bias holds across the full fine-tuning checkpoint ladder as clean success rises from $60\%$ to $74\%$ (so it is not an undertraining artifact), replicates across two training seeds, and replicates---more strongly---on a second LIBERO suite, LIBERO-object (\cref{sec:emergence,sec:generalization}).
\end{itemize}

\cref{fig:cloud} sets up the picture: under a $50$\,mm displacement the reach endpoints of a competent policy cluster at the canonical location $\Cv$, not the displaced object $\Tv$ (\cref{sec:competent}); the object-paste counterfactual (\cref{sec:sufficiency}) then shows this reach does not move \emph{even when the object's appearance does}. The remainder of the paper formalizes the diagnostic (\cref{sec:protocol}), quantifies the memory trap on a competent policy (\cref{sec:competent}), runs the causal counterfactual that is our main result (\cref{sec:sufficiency}), characterizes the bias's persistence across competence (\cref{sec:emergence}), and discusses implications and scope (\cref{sec:discussion,sec:limitations}).
